%! AP Statistics — Unit 2, Lesson 5
\documentclass[10pt]{article}
\usepackage[margin=0.6in]{geometry}
\usepackage{xcolor}
\usepackage[most]{tcolorbox}
\usepackage{enumitem}
\usepackage{multicol}
\usepackage{amsmath,amssymb}
\usepackage{array}
\usepackage{tabularx}
\tcbuselibrary{breakable}

\definecolor{sky}{HTML}{EAF3FB}
\definecolor{navy}{HTML}{1F3A5F}
\definecolor{redbox}{HTML}{FCECEC}
\definecolor{greenbox}{HTML}{EDF8F0}
\definecolor{goldbox}{HTML}{FFF7E6}
\newtcolorbox{skillbox}[2][]{
    colback=#2, colframe=navy, boxrule=0.8pt, arc=2mm,
    left=2mm,right=2mm,top=1.5mm,bottom=1.5mm,
    fonttitle=\bfseries, title={#1}, halign title=center, breakable
}
\setlist[itemize]{leftmargin=1.2em,itemsep=0.35em,topsep=0.35em}
\setlist[enumerate]{leftmargin=1.4em,itemsep=0.35em,topsep=0.35em}
\newcolumntype{C}[1]{>{\centering\arraybackslash}p{#1}}
\newcolumntype{Y}{>{\centering\arraybackslash}X}

\newcommand{\CourseName}{AP Statistics}
\newcommand{\SchoolYear}{2026--2027}
\newcommand{\MeetingLength}{55 minutes}
\newcommand{\UnitNumberName}{Unit 2: Exploring Two-Variable Data \quad}
\newcommand{\LessonNumberName}{Lesson 2.5: Correlation}

\begin{document}
\begin{center}
    {\Large\bfseries \CourseName: \SchoolYear\ (\MeetingLength) \\
    {\small\bfseries \UnitNumberName \LessonNumberName}}
\end{center}

\begin{tcolorbox}[colback=sky,colframe=navy,boxrule=0.9pt,arc=2mm,
    left=2mm,right=2mm,top=1.5mm,bottom=1.5mm]
    \textbf{Primary Objective:} Students will calculate and interpret the correlation coefficient
    $r$ as a measure of the strength and direction of a linear association between two quantitative
    variables. Students will understand the properties of $r$ (including its sensitivity to
    outliers) and recognize that correlation does not imply causation (Big Idea: DAT; AP Skills 2 \& 4).
\end{tcolorbox}

\begin{skillbox}[Priority Ideas \& Skills]{goldbox}
\begin{minipage}[t]{0.48\linewidth}
    \textbf{AP Skills 2 \& 4 --- Data Analysis \& Statistical Argumentation}
    \begin{itemize}
        \item Interpret $r$ in context: sign (direction), magnitude (strength).
        \item Recognize the properties of $r$: $-1 \leq r \leq 1$; $r = \pm 1$ means perfect linear; $r$ is sensitive to outliers; $r$ has no units.
        \item Explain why a high $|r|$ value does not establish causation.
    \end{itemize}
\end{minipage}
\hfill
\begin{minipage}[t]{0.48\linewidth}
    \textbf{Key Understandings}
    \begin{itemize}
        \item $r$ measures only \textit{linear} association; a curved relationship can have $r \approx 0$ yet be very strong.
        \item $r$ is not resistant: one influential point can dramatically change its value.
        \item $r^2$ (the coefficient of determination) tells what proportion of variation in $y$ is explained by the linear model --- preview for Lesson 2.6.
    \end{itemize}
\end{minipage}
\end{skillbox}

\begin{skillbox}[{Vocabulary, Concepts, \& Theorems}]{greenbox}
\begin{minipage}[t]{0.48\linewidth}
    \begin{itemize}
        \item \textbf{Correlation coefficient} $r$:
        \[r = \frac{1}{n-1}\sum\left(\frac{x_i - \bar{x}}{s_x}\right)\left(\frac{y_i - \bar{y}}{s_y}\right)\]
        measures linear association; $-1 \leq r \leq 1$
        \item \textbf{$r > 0$}: positive linear association; \textbf{$r < 0$}: negative linear association; \textbf{$r = 0$}: no linear association
        \item \textbf{$|r|$ close to 1}: strong; close to 0: weak (rough guides: $|r| \geq 0.8$ strong, $0.5$--$0.8$ moderate, $< 0.5$ weak)
    \end{itemize}
\end{minipage}
\hfill
\begin{minipage}[t]{0.48\linewidth}
    \begin{itemize}
        \item \textbf{$r^2$ (coefficient of determination)}: proportion of variation in $y$ explained by the linear relationship with $x$; $0 \leq r^2 \leq 1$
        \item \textbf{Properties of $r$}: unitless; symmetric ($r_{xy} = r_{yx}$); not resistant to outliers; only meaningful for linear relationships
        \item \textbf{Correlation $\neq$ causation}: a high $|r|$ can result from a confounding variable or coincidence
    \end{itemize}
\end{minipage}
\end{skillbox}

\begin{skillbox}[Activate Prior Knowledge \& Spiral Review]{sky}
\begin{minipage}[t]{0.48\linewidth}
    \textbf{Quick Review (3 min)}\\
    Display the scatterplots from Lesson 2.4 and students' DOFS descriptions. Ask: ``You used words like `strong' and `moderate.' Could we replace those words with a single number?'' That number is $r$.
\end{minipage}
\hfill
\begin{minipage}[t]{0.48\linewidth}
    \textbf{Spiral Connections}
    \begin{itemize}
        \item DOFS from Lesson 2.4 gives qualitative description; $r$ gives a quantitative one.
        \item Lesson 2.8: $r^2$ is reported alongside regression output.
        \item Unit 9: inference for the slope of a regression model also tests whether $r$ is significantly different from 0.
    \end{itemize}
\end{minipage}
\end{skillbox}

\begin{skillbox}[Hook \& Lesson]{sky}
\begin{minipage}[t]{0.48\linewidth}
    \textbf{Hook (5 min)}\\
    Show 6 scatterplots with different $r$ values ($r = -0.95, -0.5, 0, 0.5, 0.85, 1.0$) displayed out of order. Students match each plot to an $r$ value by intuition. Debrief: correlation is a formalization of the visual impression of linear tightness.
\end{minipage}
\hfill
\begin{minipage}[t]{0.48\linewidth}
    \textbf{Lesson Flow (15 min)}
    \begin{enumerate}
        \item Introduce the formula for $r$ conceptually (products of standardized values).
        \item Walk through a by-hand computation for 4 data points.
        \item Discuss the properties of $r$ one by one, with a counterexample for each potential misconception.
        \item Demonstrate with technology (calculator or Desmos correlation feature).
        \item Introduce $r^2$ briefly as ``proportion of variation explained.''
    \end{enumerate}
\end{minipage}
\end{skillbox}

\begin{skillbox}[Explicit Instruction \& Active Monitoring]{sky}
\begin{minipage}[t]{0.48\linewidth}
    \textbf{Direct Instruction (10 min)}
    \begin{itemize}
        \item Interpretation template: ``The correlation between [x] and [y] is $r = [value]$, indicating a [strong/moderate/weak] [positive/negative] linear association.''
        \item Demonstrate the outlier sensitivity of $r$: show a scatterplot, compute $r$, remove one outlier, recompute $r$ --- dramatic change.
        \item Show a curved relationship with $r \approx 0$ to illustrate that $r = 0$ does NOT mean no relationship.
    \end{itemize}
\end{minipage}
\hfill
\begin{minipage}[t]{0.48\linewidth}
    \textbf{Active Monitoring}
    \begin{itemize}
        \item Common misconception: $r = 0.9$ means 90\% of the data is linear. Correct: it describes strength; $r^2 = 0.81$ means 81\% of variation is explained.
        \item Watch for: students switching $r$ and $r^2$ in interpretations.
        \item Probe: ``If I switch $x$ and $y$, what happens to $r$?'' (Nothing --- $r$ is symmetric.)
    \end{itemize}
\end{minipage}
\end{skillbox}

\begin{skillbox}[Group Work \& Differentiation]{redbox}
\begin{minipage}[t]{0.48\linewidth}
    \textbf{Group Activity (8--10 min)}\\
    Each group receives a dataset and computes $r$ by hand (or with a calculator). Groups then:
    \begin{enumerate}
        \item Interpret $r$ fully in context.
        \item Compute $r^2$ and write a sentence interpreting it.
        \item Remove one apparent outlier, recompute $r$, and describe the change.
        \item Discuss: does the correlation establish a causal relationship?
    \end{enumerate}
\end{minipage}
\hfill
\begin{minipage}[t]{0.48\linewidth}
    \textbf{Differentiation}
    \begin{itemize}
        \item \textit{Support}: Provide a formula sheet and a worked example of the by-hand calculation; students replicate with new data.
        \item \textit{Extension}: Use an online correlation guessing game (e.g., \texttt{guessthecorrelation.com}) to build intuition for estimating $r$ from scatterplots.
        \item \textit{AP Connection}: Practice interpreting $r^2$ --- the most commonly misinterpreted statistic in AP Statistics FRQs.
    \end{itemize}
\end{minipage}
\end{skillbox}

\begin{skillbox}[Individual Work \& Assessment]{redbox}
\begin{minipage}[t]{0.48\linewidth}
    \textbf{Exit Ticket (5 min)}\\
    Given: $r = -0.74$ between daily temperature and hot chocolate sales. Students: (1) describe the association, (2) compute and interpret $r^2$, (3) identify one possible confounding variable that could explain the association without direct causation.
\end{minipage}
\hfill
\begin{minipage}[t]{0.48\linewidth}
    \textbf{AP-Style Check}\\
    \textit{``The correlation between study hours and exam score is $r = 0.88$. A student concludes that studying causes higher scores. Evaluate this conclusion.''}\\[4pt]
    Assess: student acknowledges the strong linear association, but correctly states this is an observational study and association does not establish causation (lurking variables possible).
\end{minipage}
\end{skillbox}

\begin{skillbox}[Reinforcement \& Extension]{goldbox}
\begin{minipage}[t]{0.48\linewidth}
    \textbf{Homework / Practice}
    \begin{itemize}
        \item For 3 provided datasets, compute $r$ and $r^2$, interpret both, and state whether causation can be concluded.
        \item AP textbook: multiple-choice on properties and interpretation of $r$.
    \end{itemize}
\end{minipage}
\hfill
\begin{minipage}[t]{0.48\linewidth}
    \textbf{Extension / Preview}
    \begin{itemize}
        \item \textit{Extension}: Visit \texttt{tylervigen.com/spurious-correlations} --- find two correlations that are clearly not causal. Write a brief explanation of the likely lurking variable.
        \item \textit{Preview}: Lesson 2.6 --- knowing the strength of the linear association, we now want to \textit{model} it with a line. The least-squares regression line uses the same $r$, $\bar{x}$, $\bar{y}$, $s_x$, and $s_y$ values.
    \end{itemize}
\end{minipage}
\end{skillbox}

\end{document}
